\clearpage
\fancyhead[L]{\small CONNOISSANCE DES TEMS.}
\hypertarget{native-522-lower}{}\label{native-522-lower}
\begin{center}\small Göttingische gelehrte Anzeigen. Stück 95. Seite 939..943. 1809. Juni 17.\\[3pt]\rule{.72\linewidth}{.3pt}\end{center}

\emph{Connoissance des tems ou des mouvemens célestes, à l'usage des astronomes et des navigateurs, pour l'an} 1810; \emph{publiée par le bureau des longitudes. De l'imprimerie impériale} 1808. gr. Octav. Paris.

Wir übergehen den Kalender und die übrigen stehenden Artikel, die in ihrer Einrichtung ungeändert geblieben sind, mit Stillschweigen, und bemerken nur, dass die Tafel der geopraphischen Längen und Breiten diesmal wieder einige nicht unbedeutende Bereicherungen erhalten hat, vornehmlich \SourceLeaf{523}mehrere schätzbare neue Bestimmungen aus der Insel Cypern, Arabien und dem rothen Meere. Die \emph{Additions}, um deren willen wir eigentlich diese Anzeige geben, werden, wie in den beiden vorigen Jahrgängen, mit den Beobachtungen des Hrn. \textsc{Bouvard} auf der kaiserl. Sternwarte, während des Jahres 1806, eröffnet. Sie füllen 80 Seiten: wir finden darunter, ausser der Sonne und dem Monde, auch die sämmtlichen ältern Planeten, aber diesmal gar keine von Ceres, Pallas oder Juno. Der von \textsc{Pons} im November zu Marseille entdeckte Comet wurde von \textsc{Bouvard} vom 21. November bis 19. December beobachtet. Verfinsterungen von Jupiterstrabanten ziemlich zahlreich; Sternbedeckungen zusammen fünf, und zwar blos Eintritte ($g$ Löwen den 9. Januar, $\zeta$ Krebs den 1. März, $\theta$ Ophiuchus den 1. Juni, $\xi$ Zwillinge den 8. September, 19 Fische den 20. November). — Hierauf folgen Chinesische Beobachtungen seit dem Jahre 147 vor unsrer Zeitrechnung, eingesandt von P. \textsc{Gaubil} im Jahre 1749; diese Beobachtungen beschränken sich auf nahe Zusammenkünfte und Bedeckungen des Mondes und der Planeten unter einander und mit Fixsternen, ohne nähere Umstände, und grössten Theils nur mit Angabe des Tages, ohne die Stunde: eigentlicher Nutzen wird also wenig daraus zu ziehen sein. — Eine literarische Notiz über \textsc{Hevel}'s und \textsc{Dörfel}'s Verdienste um die Theorie der Bewegung der Cometen, von \textsc{J. C. Burckhardt}. — Sechste und letzte Sammlung der Beobachtungen \textsc{Messier}'s, von 1752 bis Ende 1759: zwei Mondsfinsternisse, viele Sternbedeckungen vom Monde, noch viel zahlreichere Verfinsterungen von Jupiterstrabanten, zwei Cometen, wovon der eine, der \textsc{Halley}'sche, am 21. Januar 1759 zuerst aufgefunden wurde: auf Verlangen seines Lehrers \textsc{Delisle} musste \textsc{Messier} diese Auffindung bis Anfang April geheim halten. — Beobachtung der obern Zusammenkunft des Mercur und der untern der Venus im September und October 1807, von \textsc{Vidal} zu Mirepoix. — Beobachtung des grossen Cometen von 1807, von demselben, vom 27. November bis 4. März 1808. — Beobachtungen desselben Cometen auf der königl. Lissaboner Sternwarte von \textsc{Paul Ciera} (vom 7. October bis 29. November), und zu Bremen von \textsc{Olbers} (vom 8. October bis 14. Februar 1808). — Messung eines Erdmeridianbogens und eines auf den Meridian senkrechten Grades in Ostindien, vom Brigademajor \textsc{William Lambton}: ein Auszug aus den Memoirs von Calcutta. Die Basis hielt 40006 Engl. Fuss, die Endpunkte des gemessenen Meridianbogens waren

\begin{center}\begin{tabular}{@{}lll@{}}\emph{Paudree} & Breite & $13^\circ\,19'\,49''018$ \\ \emph{Trivandepcrum} & Breite & $11\quad44\quad52.\,59$ \\ \end{tabular}\end{center}
\noindent Man leitete daraus den Breitengrad ab zu 60495 Fathoms oder 56763 Toisen; den Längengrad fand man 61061 Fathoms oder 57294 Toisen. Aus einem wie grossen Bogen lezterer bestimmt wurde, ist nicht angegeben, daher wir über den Grad der Genauigkeit, welcher ihm beigelegt werden darf, nicht urtheilen können. Die Vergleichung beider gäbe die zu starke Abplattung $\frac{1}{205,67}$, mit der Französischen Gradmessung verglichen, gäbe der Breitengrad die Abplattung $\frac{1}{817}$. — Beobachtungen auf der Lissaboner Sternwarte, von \textsc{Paul Ciera} (ausser drei Bedeckungen des Sterns $\alpha^2$ Krebs vom Monde, 1807 Februar 20, April 16, October 14, blos Jupiterstrabanten-Verfinsterungen). — Methode zur Berechnung der Correctionen der Durchgänge am Mittagsfernrohr, von \textsc{Delambre}: ein sehr weitläuftiger Aufsatz, der indess für weniger geübte Practiker seinen Nutzen haben mag. Dass das Mittagsfernrohr blos durch Vergleichung mehrerer Durchgänge mit den Rectascensionen, ohne Zuziehung des Niveaus oder der absoluten Sternzeit, oder sonst Etwas, was sich auf das Zenith des Beobachters bezieht, nicht in den Meridian, sondern blos in einen Stundenkreis gebracht werden könne, scheint Hr. \textsc{Delambre} für eine neue Bemerkung zu halten; wir können aber nicht glauben, dass etwas so Offenbares je einem nachdenkenden Astronomen entgangen sein könne. — Die Mondfinsterniss vom 4. Januar 1806, beobachtet von \textsc{Vidal} zu Mirepoix. — Beobachtungen auf der kaiserl. Sternwarte zu Marseille, von \textsc{Thulis}, von 1796..\SourceLeaf{524}1806. — Beobachtung des grossen Cometen 1807 zu Montauban, von \textsc{Duc-la Chapelle}. — Beobachtungen zu Viviers, von \textsc{Flaugergues}, im Jahre 1807. — Die allgemeinen Aberrations- und Nutationstafeln des Hrn. \textsc{von Zach}, aus dessen bekanntem Werke abgedruckt. Eben so die allgemeinen Aberrations- und Nutationstafeln von \textsc{Gauss}, aus der Monatl. Correspondenz abgedruckt, nebst einer Entwickelung der Formeln, auf denen sie beruhen. — Geschichte der Astronomie für 1808: besteht diesmal blos in Auszügen aus neu erschienenen astronomischen Werken: \textsc{von Zach}'s \emph{Tabulae speciales}, wo Hr. \textsc{Delambre} Gelegenheit nimmt, die Formeln, auf denen seine in jenem Werke wieder benutzten Aberrationstafeln für die Planeten beruhen, mitzutheilen; \textsc{Cagnoli}'s \emph{Sternverzeichniss}; \textsc{Monteiro-da-Rocha}'s \emph{Mémoires sur l'astronomie pratique}, aus dem Portugiesischen ins Französische übersetzt; die Französische Uebersetzung von \textsc{Cagnoli}'s Trigonometrie, nach der zweiten Ausgabe; die erste Lieferung des astronomischen Theils von \textsc{Alexander von Humboldt}'s Reise; \textsc{L. W. Pfaff} \emph{de tubo Culminatorio Dorpatensi brevis narratio}, \emph{Dorpati} 1808; \emph{Efemeridi astronomiche di Milano per l'anno} 1809. Man sieht aus dieser Aufzählung, dass die \emph{Connoissance des tems}, wenn gleich dieser Band an Originalaufsätzen nicht ganz so reich ist, wie manche seiner Vorgänger, doch noch immer fortfährt, ein den Astronomen sehr schätzbares Repertorium zu sein, worin man Nachrichten über alles das Wissenswürdigste, was in der Astronomie geschieht, nicht umsonst suchen wird. — Auf diese Auszüge folgt noch ein Bericht von den Endresultaten der bis zur Insel Formentera fortgesetzten Französischen Gradmessung, wobei ein Dreieck gebraucht wurde, dessen eine Seite 82555 Toisen hielt. Der ganze Bogen von Dünkirchen bis Formentera ist dadurch bis auf $13^\circ\,12'\,22''$ angewachsen, und der Meter wird dadurch von 443,296 Linien auf 443,2958 Linien vermindert, d. i. so gut als gar nicht verändert. Den Beschluss dieses Bandes machen einige Verbesserungen zu den vom Bureau des Longitudes herausgegebenen Sonnen- und Mondtafeln, eine neu berechnete Tafel für die Zeitgleichung auf 1800, nebst den Säcularänderungen, und die von \textsc{Bouvard} im Jahr 1806 angestellten meteorologischen Beobachtungen.
